% ======================================================================
%  sec26.tex —— 第 26 节 Proj(S) 上的拟凝聚模
%  底本：原书 pp.124–132；原书标题：26. Quasi-coherent modules on Proj(S)
% ======================================================================

% ---------- 本节用到的宏（若全书导言区已定义则自动跳过） ----------
\providecommand{\QCoh}{\mathsf{QCoh}}
\providecommand{\Sym}{\operatorname{Sym}}
\providecommand{\Ann}{\operatorname{Ann}}
% 原书 26.2.2 的环境名为 Proposition-Definition，故单立一个共享计数器的环境
\makeatletter
\@ifundefined{propdef}{%
  \theoremstyle{definition}%
  \newtheorem{propdef}[theorem]{命题—定义}%
  \crefname{propdef}{命题—定义}{命题—定义}%
}{}
\makeatother

\section{$\Proj(S)$ 上的拟凝聚模}\label{sec:26}

本节固定一个 $\Z_{\ge 0}$-分次环 $S$。记 $\cO := \cO_{\Proj(S)}$。

\subsection{$\widetilde{M}$}\label{sec:26.1}

\begin{construction}\label{con:26.1.1}
设 $M$ 为 $\Z$-分次 $S$-模。考虑拟凝聚模 $\widetilde{M} \in \QCoh(\Spec(S))$ 及其限制
\[
\widetilde{M}|_{\Spec(S)^{\circ}} \in \QCoh(\Spec(S)^{\circ}).
\]
回忆典范的仿射态射 $\pi \colon \Spec(S)^{\circ} \to \Proj(S) = X$。由此得到拟凝聚模
\begin{equation}\label{eq:26.1}
\pi_*\bigl(\widetilde{M}|_{\Spec(S)^{\circ}}\bigr) \in \QCoh(\Proj(S)).
\end{equation}
对任意标准开子集 $V \subseteq \Proj(S)$，取齐次元素 $f \in S_+$ 使 $V = U_+(f)$。则
\[
\bigl(\pi_*(\widetilde{M}|_{\Spec(S)^{\circ}})\bigr)(V) \simeq \widetilde{M}(U(f)) \simeq M_f.
\]
\end{construction}

这在 $\bigl(\pi_*(\widetilde{M}|_{\Spec(S)^{\circ}})\bigr)(V)$ 上给出一个 $\Z$-分次。可以验证该分次不依赖 $f$ 的选取，且与限制映射相容。于是 (26.1) 上有一个 $\Z$-分次。我们滥用记号，定义
\[
\widetilde{M} := \bigl(\pi_*(\widetilde{M}|_{\Spec(S)^{\circ}})\bigr)_0 \in \QCoh(\Proj(S))
\]
为 (26.1) 的零次部分，并称之为 $\Proj(S)$ 上\emph{与分次模 $M$ 相伴的拟凝聚模}。

\begin{remark}\label{rem:26.1.2}
由构造，在任意标准开子集 $U_+(f) \subseteq \Proj(S)$ 上，有
\[
\widetilde{M}|_{U_+(f)} \simeq \widetilde{M_{(f)}} \in \QCoh(U_+(f)),
\]
其中 $M_{(f)} := (M_f)_0$ 是 $S_{(f)}$-模。

对 $f \in S_n$，有
\[
M_{(f)} \simeq \Bigl(\bigoplus_{d \in \Z} M_{dn}\Bigr)/(f-1).
\]
\end{remark}

\begin{example}\label{ex:26.1.3}
取 $M = S$，则 $\widetilde{M} \simeq \cO$。
\end{example}

注意 $\Z$-分次 $S$-模构成阿贝尔范畴。上述构造对 $M$ 具有函子性，即得到函子
\begin{equation}\label{eq:26.2}
S\text{-}\mathsf{mod}_{\mathrm{gr}} \to \QCoh(\Proj(S)),
\end{equation}
它把 $\Z$-分次 $S$-模 $M$ 送到拟凝聚 $\cO$-模 $\widetilde{M}$。

下述结果由注记 \ref{rem:26.1.2} 得到。

\begin{lemma}\label{lem:26.1.4}
函子 (26.2) 正合，且与（小）余极限交换。
\end{lemma}

\begin{warning}\label{warn:26.1.5}
函子 (26.2) 不是等价。事实上，它甚至不保守。
\end{warning}

\begin{definition}\label{def:26.1.6}
称 $\Z$-分次 $S$-模 $M$ \emph{最终为零}，若当 $n \gg 0$ 时 $M_n = 0$。
\end{definition}

\begin{lemma}\label{lem:26.1.7}
若 $M$ 最终为零，则 $\widetilde{M} = 0$。
\end{lemma}

\begin{proof}
对任意齐次 $f \in S_+$，有 $M_{(f)} = 0$。于是 $\widetilde{M}|_{U_+(f)} = 0$，从而 $\widetilde{M} = 0$。
\end{proof}

\begin{notation}\label{not:26.1.8}
设 $M$ 为 $\Z$-分次 $S$-模，$n$ 为整数。记 $M(n)$ 为由下式定义的 $\Z$-分次 $S$-模：
\[
M(n)_m := M_{m+n}.
\]
\end{notation}

\begin{remark}\label{rem:26.1.9}
由构造，有
\[
\pi_*\bigl(\widetilde{M}|_{\Spec(S)^{\circ}}\bigr) \simeq \bigoplus_{n \in \Z} \widetilde{M(n)}
\]
（作为 $\Z$-分次拟凝聚 $\cO$-模）。
\end{remark}

\begin{lemma}\label{lem:26.1.10}
设 $U_+(f) \subseteq \Proj(S)$ 为标准开子集，$M$ 为 $\Z$-分次 $S$-模，且 $M \to M_f$ 是双射。则 $\widetilde{M}$ 是沿开浸入 $U_+(f) \to \Proj(S)$ 的 $\ast$-延拓。
\end{lemma}

\begin{proof}
该假设蕴涵 $M$ 是 $S_f$-模。于是对象 $\widetilde{M} \in \QCoh(\Spec(S)^{\circ})$ 是沿开浸入 $U(f) \to \Spec(S)^{\circ}$ 的 $\ast$-延拓。回忆 $U(f)$ 在 $\pi \colon \Spec(S)^{\circ} \to \Proj(S)$ 下的像是 $U_+(f)$，故 $\widetilde{M} \in \QCoh(\Proj(S))$ 是沿开浸入 $U_+(f) \to \Proj(S)$ 的 $\ast$-延拓。
\end{proof}

\subsection{重言线丛}\label{sec:26.2}

本小节假设理想 $S_+$ 由 $S_1$ 中元素生成。

\begin{warning}\label{warn:26.2.1}
若 $S_+$ 不由 $S_1$ 生成，本小节全部结果都可能失效。
\end{warning}

\begin{propdef}\label{prop:26.2.2}
对任意 $n \in \Z$，$\cO$-模
\[
\cO(n) := \widetilde{S(n)}
\]
可逆。我们称 $\cO(1)$ 为 $\Proj(S)$ 上的\emph{重言线丛}。
\end{propdef}

\begin{proof}
关于 $S$ 的假设蕴涵 $\Proj(S)$ 可被 $f \in S_1$ 对应的 $U_+(f)$ 覆盖。故只需证 $\widetilde{S(n)}|_{U_+(f)}$ 局部自由、秩为 $1$。按定义，它是对应于 $S_{(f)}$-模 $S(n)_{(f)}$ 的拟凝聚 $\cO_{U_+(f)}$-模。

因 $f$ 的次数为 $1$，对任意 $d \in \Z$，映射 $f \colon (S_f)_d \to (S_f)_{d+1}$ 与 $f^{-1} \colon (S_f)_{d+1} \to (S_f)_d$ 互为逆。于是
\[
S(n)_{(f)} \simeq (S_f)_n \simeq (S_f)_0 \simeq S_{(f)}.
\]
这说明 $\widetilde{S(n)}|_{U_+(f)} \simeq \cO_{U_+(f)}$。
\end{proof}

\begin{yizhu}
证明中「$f \in S_1$」这一条不可省：正因为 $\deg f = 1$，$S_f$ 的 $\Z$-分次才由 $f$ 的倍数平移给出，$S(n)_{(f)} \simeq (S_f)_0$ 才成立。若 $S_+$ 不由 $S_1$ 生成，可取 $f$ 的次数为正整数 $d > 1$，此时 $\cO(n)$ 一般不可逆，乘法映射 $\cO(m) \otimes \cO(n) \to \cO(m+n)$ 也一般不同构；能断言的是 $\cO_X(dn)$ 在开集 $W_d = \bigcup_{f \in S_d} D_+(f)$ 上可逆。
[来源：Stacks Project, Tag 01MM（Definition 27.10.1 及紧随其后的讨论）与 Tag 01MU（Lemma 27.10.4），\url{https://stacks.math.columbia.edu/tag/01MM}、\url{https://stacks.math.columbia.edu/tag/01MU}（2026-11-17 访问）。]
\end{yizhu}

\begin{corollary}\label{cor:26.2.3}
对任意 $n \in \Z$，有自同构
\[
\QCoh(\Proj(S)) \to \QCoh(\Proj(S)), \qquad \cF \mapsto \cF(n) := \cF \otimes_\cO \cO(n).
\]
\end{corollary}

\begin{exercise}\label{exc:26.2.4}
对任意 $\Z$-分次 $S$-模 $M$，有
\[
\pi^*(\widetilde{M}) \simeq \widetilde{M}|_{\Spec(S)^{\circ}}
\]
（作为 $\QCoh(\Spec(S)^{\circ})$ 中的对象）。
\end{exercise}

\begin{construction}\label{con:26.2.5}
设 $M$、$N$ 为 $\Z$-分次 $S$-模。回忆 $\pi \colon \Spec(S)^{\circ} \to \Proj(S)$ 是仿射的。由习题 20.2.3，有同构
\begin{equation}\label{eq:26.3}
\begin{aligned}
\pi_*\bigl(\widetilde{M}|_{\Spec(S)^{\circ}}\bigr) \otimes_{\pi_*(\cO_{\Spec(S)^{\circ}})} \pi_*\bigl(\widetilde{N}|_{\Spec(S)^{\circ}}\bigr)
&\simeq \pi_*\bigl(\widetilde{M}|_{\Spec(S)^{\circ}} \otimes_{\cO_{\Spec(S)^{\circ}}} \widetilde{N}|_{\Spec(S)^{\circ}}\bigr)\\
&\simeq \pi_*\bigl(\widetilde{M \otimes_S N}|_{\Spec(S)^{\circ}}\bigr)
\end{aligned}
\end{equation}
（在 $\QCoh(\Proj(S))$ 中）。而且可以验证这些同构与这些对象上自然的 $\Z$-分次相容。于是得到 $\QCoh(\Proj(S))$ 中的态射
\begin{equation}\label{eq:26.4}
\begin{aligned}
\pi_*\bigl(\widetilde{M}|_{\Spec(S)^{\circ}}\bigr)_0 \otimes_{\pi_*(\cO_{\Spec(S)^{\circ}})_0} \pi_*\bigl(\widetilde{N}|_{\Spec(S)^{\circ}}\bigr)_0
&\to \pi_*\bigl(\widetilde{M}|_{\Spec(S)^{\circ}}\bigr) \otimes_{\pi_*(\cO_{\Spec(S)^{\circ}})} \pi_*\bigl(\widetilde{N}|_{\Spec(S)^{\circ}}\bigr)\\
&\simeq \pi_*\bigl(\widetilde{M \otimes_S N}|_{\Spec(S)^{\circ}}\bigr)_0.
\end{aligned}
\end{equation}
换言之，有典范态射
\begin{equation}\label{eq:26.5}
\widetilde{M} \otimes_\cO \widetilde{N} \to \widetilde{M \otimes_S N}
\end{equation}
（在 $\QCoh(\Proj(S))$ 中）。
\end{construction}

\begin{lemma}\label{lem:26.2.6}
典范态射 (26.5) 是同构。
\end{lemma}

\begin{proof}
可限制到 $f \in S_1$ 对应的标准开子集 $U_+(f)$。此时 (26.5) 是对应于
\[
M_{(f)} \otimes_{S_{(f)}} N_{(f)} \to (M \otimes_S N)_{(f)}
\]
的态射，容易验证它是同构。
\end{proof}

\begin{remark}\label{rem:26.2.7}
事实上可以证明
\[
S\text{-}\mathsf{mod}_{\mathrm{gr}} \to \QCoh(\Proj(S)), \quad M \mapsto \widetilde{M}
\]
是对称幺半函子。
\end{remark}

\begin{corollary}\label{cor:26.2.8}
有
\[
\cO(m) \otimes_\cO \cO(n) \simeq \cO(m+n).
\]
特别地，对 $\cF \in \QCoh(\Proj(S))$，有
\[
(\cF(m))(n) \simeq \cF(m+n).
\]
\end{corollary}

\begin{corollary}\label{cor:26.2.9}
设 $M$ 为 $\Z$-分次 $S$-模。则
\[
\widetilde{M}(n) \simeq \widetilde{M(n)}.
\]
\end{corollary}

\begin{definition}\label{def:26.2.10}
设 $M$ 为 $\Z$-分次 $S$-模。称 $M$ \emph{（分次）有限表现}，若存在 $\Z$-分次 $S$-模的正合列
\[
\bigoplus_{i=1}^{m} S(d_i) \to \bigoplus_{j=1}^{n} S(e_j) \to M \to 0.
\]
\end{definition}

\begin{exercise}\label{exc:26.2.11}
设 $M$ 为有限表现的 $\Z$-分次 $S$-模。证明：
\begin{enumerate}[label=(\arabic*),leftmargin=2.4em,itemsep=0.25ex]
  \item $\cO$-模 $\widetilde{M}$ 有限表现。
  \item 对任意 $\Z$-分次 $S$-模 $N$，有
  \[
  \uHom_{\cO}\bigl(\widetilde{M}, \widetilde{N}\bigr) \simeq \widetilde{\uHom^{\mathrm{gr}}_{S}(M, N)},
  \]
  其中 $\uHom^{\mathrm{gr}}_{S}(M, N)$ 是使
  \[
  \uHom^{\mathrm{gr}}_{S}(M, N)_d := \Hom_{S\text{-}\mathsf{mod}_{\mathrm{gr}}}\bigl(M, N(d)\bigr)
  \]
  的 $\Z$-分次 $S$-模。
\end{enumerate}
\end{exercise}

\subsection{扭转的整体截面}\label{sec:26.3}

本小节假设理想 $S_+$ 由 $S_1$ 中元素生成。

\begin{construction}\label{con:26.3.1}
注意 $\bigoplus_{n \in \Z} \cO(n)$ 自然地是 $\Proj(S)$ 上 $\Z$-分次的交换环层。于是
\[
\Gamma(\cO)_{\mathrm{gr}} := \bigoplus_{n \in \Z} \Gamma\bigl(\Proj(S), \cO(n)\bigr)
\]
自然地是 $\Z$-分次交换环。

对任意 $\cO$-模 $\cF$，$\bigoplus_{n \in \Z} \cF(n)$ 自然地是 $\bigoplus_{n \in \Z} \cO(n)$ 上的 $\Z$-分次模。于是
\[
\Gamma(\cF)_{\mathrm{gr}} := \bigoplus_{n \in \Z} \Gamma\bigl(\Proj(S), \cF(n)\bigr)
\]
自然地是 $\Gamma(\cO)_{\mathrm{gr}}$ 上的 $\Z$-分次模。

上述构造对 $\cF$ 具有函子性。换言之，有函子
\[
\Gamma(-)_{\mathrm{gr}} \colon \cO\text{-}\mathsf{mod} \to \Gamma(\cO)_{\mathrm{gr}}\text{-}\mathsf{mod}_{\mathrm{gr}}.
\]
\end{construction}

\begin{construction}\label{con:26.3.2}
设 $M$ 为 $\Z$-分次 $S$-模。

对任意标准开子集 $V \subseteq \Proj(S)$，由注记 \ref{rem:26.1.9}，有阿贝尔群之间的典范态射
\begin{equation}\label{eq:26.6}
M \to \Gamma\bigl(\pi^{-1}(V), \widetilde{M}\bigr) \simeq \Gamma\Bigl(V, \bigoplus_{n \in \Z} \widetilde{M(n)}\Bigr) \simeq \bigoplus_{n \in \Z} \Gamma(V, \widetilde{M(n)}).
\end{equation}
可取齐次元素 $f \in S_+$ 使 $V = U_+(f)$。此时 (26.6) 可等同于典范态射 $M \to M_f$。由此 (26.6) 与两侧的 $\Z$-分次相容。

由构造，(26.6) 与限制映射相容。因此得到 $\Z$-分次态射
\begin{equation}\label{eq:26.7}
M \to \bigoplus_{n \in \Z} \Gamma\bigl(\Proj(S), \widetilde{M(n)}\bigr) = \Gamma(\widetilde{M})_{\mathrm{gr}}.
\end{equation}
在 $M = S$ 的特殊情形，可以验证所得映射
\begin{equation}\label{eq:26.8}
S \to \Gamma(\cO)_{\mathrm{gr}}
\end{equation}
是 $\Z$-分次交换环之间的同态。而且 (26.7) 与 (26.8) 在两侧的作用相容。换言之，(26.7) 是 $\Z$-分次 $S$-模之间的态射，其中 $S$ 经由同态 (26.8) 作用在 $\Gamma(\widetilde{M})_{\mathrm{gr}}$ 上。

态射 (26.7) 对 $M$ 具有函子性。换言之，我们得到自然变换
\begin{equation}\label{eq:26.9}
\id \to \Gamma(\widetilde{-})_{\mathrm{gr}},
\end{equation}
它位于 $S\text{-}\mathsf{mod}_{\mathrm{gr}}$ 的两个自函子之间。
\end{construction}

\begin{proposition}\label{prop:26.3.3}
自然变换 (26.9) 诱导伴随对
\[
\begin{array}{ccc}
S\text{-}\mathsf{mod}_{\mathrm{gr}} & \leftrightarrows & \cO\text{-}\mathsf{mod}\\[0.4ex]
M & \mapsto & \widetilde{M}\\
\Gamma(\cF)_{\mathrm{gr}} & \mapsfrom & \cF.
\end{array}
\]
\end{proposition}

\begin{proof}
对任意 $\Z$-分次 $S$-模 $M$ 与任意 $\cO$-模 $\cF$，下述数据等价：
\begin{enumerate}[label=(\roman*),leftmargin=2.6em,itemsep=0.25ex]
  \item 一个 $\Z$-分次 $S$-模之间的态射
  \[
  M \to \Gamma(\cF)_{\mathrm{gr}}.
  \]
  \item 对任意标准开子集 $V \subseteq \Proj(S)$，一个 $\Z$-分次 $S$-模之间的态射
  \[
  M \to \bigoplus_{n \in \Z} \Gamma(V, \cF(n)),
  \]
  且这些态射对 $V$ 与限制映射相容。
\end{enumerate}
由于 $\Proj(S)$ 被 $f \in S_1$ 对应的 $U_+(f)$ 覆盖，在 (ii) 中可设 $V = U_+(f)$，其中 $f \in S_1$。

展开定义，$S$ 经由同态
\[
S \to \bigoplus_{n \in \Z} \Gamma(\Proj(S), \cO(n)) \xrightarrow{\ \mathrm{res}\ } \bigoplus_{n \in \Z} \Gamma(V, \cO(n))
\]
作用在 $\bigoplus_{n \in \Z} \Gamma(V, \cF(n))$ 上。于是 (ii) 等价于下述数据：
\begin{enumerate}[label=(\roman*),start=3,leftmargin=2.6em,itemsep=0.25ex]
  \item 对任意标准开子集 $V \subseteq \Proj(S)$，一个 $\Z$-分次 $\bigoplus_{n \in \Z} \Gamma(V, \cF(n))$-模之间的态射
  \[
  \Bigl(\bigoplus_{n \in \Z} \Gamma(V, \cO(n))\Bigr) \otimes_S M \to \bigoplus_{n \in \Z} \Gamma(V, \cF(n)),
  \]
  且这些态射对 $V$ 与限制映射相容。
\end{enumerate}
我们断言有典范同构
\[
\Bigl(\bigoplus_{n \in \Z} \Gamma(V, \cO(n))\Bigr) \otimes_S M \simeq \bigoplus_{n \in \Z} \Gamma(V, \widetilde{M(n)}).
\]
确实，取 $f \in S_1$ 使 $V = U_+(f)$，该同构由
\[
\Bigl(\bigoplus_{n \in \Z} S(n)_{(f)}\Bigr) \otimes_S M \simeq S_f \otimes_S M \simeq M_f \simeq \bigoplus_{n \in \Z} M(n)_{(f)}
\]
给出。由于 $V$ 仿射，这说明 (iii) 等价于下述数据：
\begin{enumerate}[label=(\roman*),start=4,leftmargin=2.6em,itemsep=0.25ex]
  \item 对任意标准开子集 $V \subseteq \Proj(S)$，一个 $\Z$-分次 $\bigoplus_{n \in \Z} \cO(n)|_V$-模之间的态射
  \[
  \bigoplus_{n \in \Z} \widetilde{M(n)}|_V \to \bigoplus_{n \in \Z} \cF(n)|_V,
  \]
  且这些态射对 $V$ 与限制映射相容。
\end{enumerate}
这又等价于
\begin{enumerate}[label=(\roman*),start=5,leftmargin=2.6em,itemsep=0.25ex]
  \item 一个 $\Z$-分次 $\bigoplus_{n \in \Z} \cO(n)$-模之间的态射
  \[
  \bigoplus_{n \in \Z} \widetilde{M(n)} \to \bigoplus_{n \in \Z} \cF(n).
  \]
\end{enumerate}
因 $\cO(1)$ 可逆，这等价于
\begin{enumerate}[label=(\roman*),start=6,leftmargin=2.6em,itemsep=0.25ex]
  \item 一个 $\cO$-模之间的态射
  \[
  \widetilde{M} \to \cF.
  \]
\end{enumerate}
\end{proof}

\subsection{有限型情形}\label{sec:26.4}

\begin{definition}\label{def:26.4.1}
设 $M$ 为 $\Z$-分次 $S$-模。称 $M$ \emph{最终有限型}，若存在整数 $N$ 使
\[
\bigoplus_{n \ge N} M_n
\]
是有限型 $S$-模。
\end{definition}

\begin{lemma}\label{lem:26.4.2}
设 $S_+$ 由有限多个元素生成，$M$ 为最终有限型的 $\Z$-分次 $S$-模。则 $\widetilde{M}$ 有限型。
\end{lemma}

\begin{proof}
由引理 \ref{lem:26.1.7} 与引理 \ref{lem:26.1.4}，可用 $\bigoplus_{n \ge N} M_n$（$N$ 充分大）替换 $M$，从而设 $M$ 有限型。

只需证对任意标准开子集 $U_+(f) \subseteq \Proj(S)$，$\widetilde{M}|_{U_+(f)}$ 有限型。为此只需证 $M_{(f)}$ 是有限生成 $S_{(f)}$-模。这是交换代数中的简单习题。
\end{proof}

\begin{exercise}\label{exc:26.4.3}
设 $S_+$ 由 $S_1$ 中有限多个元素生成，$M$ 为最终有限型的 $\Z$-分次 $S$-模。则 $\widetilde{M} = 0$ 当且仅当 $M$ 最终为零。
\end{exercise}

\begin{warning}\label{warn:26.4.4}
若去掉「$S_+$ 由 $S_1$ 生成」的假设，习题 \ref{exc:26.4.3} 不成立。
\end{warning}

\begin{example}\label{ex:26.4.5}
设 $S$ 是只集中在偶数次齐次部分的 $\Z_{\ge 0}$-分次环。则对任意只集中在奇数次齐次部分的 $\Z$-分次 $S$-模 $M$，在 $\QCoh(\Proj(S))$ 中有 $\widetilde{M} \simeq 0$。
\end{example}

\begin{yizhu}
理由在于注记 \ref{rem:26.1.2} 的公式：$S_+$ 中齐次元素的次数皆为偶数，故对 $f \in S_+$ 与 $d \in \Z$，$M_{d\deg f}$ 只在 $d\deg f$ 为奇数时非零，而 $M_{(f)}$ 只取 $d = 0$ 一项，于是 $M_{(f)} = 0$。所有标准开集上都为零，故层为零，但 $M$ 本身非零——这正是警告 \ref{warn:26.1.5} 所谓函子不保守。
[译者按（上下文推断，无外部依据）：原文只给结论，未给理由。]
\end{yizhu}

\begin{theorem}\label{thm:26.4.6}
设 $S_+$ 由 $S_1$ 中有限多个元素生成。则函子
\[
\Gamma(-)_{\mathrm{gr}} \colon \QCoh(\Proj(S)) \to S\text{-}\mathsf{mod}_{\mathrm{gr}}
\]
全忠实。
\end{theorem}

\begin{proof}
由命题 \ref{prop:26.3.3}，只需证对任意 $\cF \in \QCoh(\Proj(S))$，典范自然变换
\begin{equation}\label{eq:26.10}
\Gamma(\widetilde{-})_{\mathrm{gr}} \to \id
\end{equation}
把 $\cF$ 送到同构。

事实上，由引理 \ref{lem:26.1.4}，函子 $(-)\widetilde{\ }$ 正合；由命题 \ref{prop:26.3.3}，作为右伴随，函子 $\Gamma(-)_{\mathrm{gr}}$ 左正合。故合成 $\Gamma(\widetilde{-})_{\mathrm{gr}}$ 左正合。

由假设，可取有限多个元素 $f_i \in S_1$（$i \in I$）使
\[
\Proj(S) = \bigcup_{i \in I} U_+(f_i).
\]
记 $U_i := U_+(f_i)$，并记 $g_i \colon U_i \to \Proj(S)$ 与 $g_{ij} \colon U_i \cap U_j \to \Proj(S)$ 为开浸入。对任意 $\cO$-模 $\cF$，有对 $\cF$ 具有函子性的正合列
\[
0 \to \cF \to \bigoplus_{i \in I} g_{i,*}(\cF|_{U_i}) \to \bigoplus_{(i,j) \in I^2} g_{ij,*}(\cF|_{U_i \cap U_j}).
\]
故只需证 (26.10) 把每个 $g_{i,*}(\cF|_{U_i})$ 与每个 $g_{ij,*}(\cF|_{U_i \cap U_j})$ 送到同构。

设 $f \in S_1$，$U := U_+(f)$，$g \colon U \to \Proj(S)$ 为开浸入。只需证：对任意 $\cM \in \QCoh(U_+(f))$，自然变换 (26.10) 把 $\cM$ 送到同构。回忆 $U \simeq \Spec(S_{(f)})$ 仿射。令
\[
M = \Gamma(U, \cM)
\]
为对应的 $S_{(f)}$-模。

先证 $(\Gamma(g_*\cM)_{\mathrm{gr}})\widetilde{\ }$ 是沿 $g \colon U \to \Proj(S)$ 的 $\ast$-延拓。按定义有
\begin{equation}\label{eq:26.11}
\begin{aligned}
\Gamma(g_*\cM)_{\mathrm{gr}} &:= \bigoplus_{n \in \Z} \Gamma\bigl(\Proj(S), (g_*\cM)(n)\bigr) \simeq \bigoplus_{n \in \Z} \Gamma\bigl(U, \cM \otimes_{\cO_U} \cO(n)|_U\bigr)\\
&\simeq \bigoplus_{n \in \Z} M \otimes_{S_{(f)}} S(n)_{(f)} \simeq M \otimes_{S_{(f)}} S_f.
\end{aligned}
\end{equation}
特别地，$\Gamma(g_*\cM)_{\mathrm{gr}}$ 是 $S_f$-模。故由引理 \ref{lem:26.1.10}，$(\Gamma(g_*\cM)_{\mathrm{gr}})\widetilde{\ }$ 是沿 $g \colon U \to \Proj(S)$ 的 $\ast$-延拓。

剩下要证典范态射
\[
(\Gamma(g_*\cM)_{\mathrm{gr}})\widetilde{\ }|_U \to \cM
\]
是同构。因 $U \simeq \Spec(S_{(f)})$ 仿射，只需证对应的 $S_{(f)}$-模之间的态射是同构。由 (26.11)，该态射就是显然的同构
\[
(M \otimes_{S_{(f)}} S_f)_0 \simeq M.
\]
\end{proof}

\begin{yizhu}
定理 \ref{thm:26.4.6} 的全部内容其实是：由 $\cF$ 造出 $M := \Gamma(\cF)_{\mathrm{gr}}$ 后，典范态射 $\widetilde{M} \to \cF$ 是同构。Stacks Project 把这一步单独提出来：对任意拟凝聚 $\cF$ 与 $M := \bigoplus_{n \in \Z} \Gamma(X, \cF(n))$ 都有典范映射 $\widetilde{M} \to \cF$，它在 $0$ 次分支上是恒等，且对 $\cF$ 具有函子性；本节所加的「$S_+$ 由 $S_1$ 中有限多个元素生成」正是使该映射可逐标准开集检验、从而成为同构的条件。
[来源：Stacks Project, Tag 0B5I（Lemma 27.10.7，the map from a graded module to a quasi-coherent module），\url{https://stacks.math.columbia.edu/tag/0B5I}（2026-11-17 访问）。]
\end{yizhu}

\begin{corollary}\label{cor:26.4.7}
设 $S_+$ 由 $S_1$ 中有限多个元素生成。则任意 $\cF \in \QCoh(\Proj(S))$ 都形如 $\widetilde{M}$。而且，若 $\cF$ 有限型，则可取 $M$ 使 $M$ 有限型。
\end{corollary}

\begin{proof}
第一个断言由定理 \ref{thm:26.4.6} 立即得到。对第二个断言，取 $\Z$-分次 $S$-模使 $\cF \simeq \widetilde{M}$。把 $M = \bigcup M_\alpha$ 写成 $\Z$-分次子模之并，使每个 $M_\alpha$ 由有限多个齐次元素生成。有 $\widetilde{M} \simeq \bigcup \widetilde{M_\alpha}$。利用 $\Proj(S)$ 拟紧与 $\widetilde{M}$ 有限型，容易看出当指标 $\alpha$ 充分大时 $\widetilde{M} = \widetilde{M_\alpha}$。
\end{proof}

\begin{proposition}\label{prop:26.4.8}
设 $S_+$ 由 $S_1$ 中有限多个元素生成。对任意有限型的 $\cF \in \QCoh(\Proj(S))$，存在整数 $N$，使对 $n \ge N$，$\cO_X$-模 $\cF(n)$ 由有限多个整体截面生成，即存在形如 $\cO^{\oplus m} \to \cF(n)$ 的满射。
\end{proposition}

\begin{proof}
由推论 \ref{cor:26.4.7}，可设 $\cF = \widetilde{M}$，其中 $M$ 是某个有限型的 $\Z$-分次 $S$-模。注意存在形如 $\bigoplus_{i \in I} S(n_i) \to M$ 的满射，其中 $|I| < \infty$。故只需对每个 $\widetilde{S(n_i)}$ 证明该断言。我们把它留作读者的练习。
\end{proof}

\begin{exercise}\label{exc:26.4.9}
设 $S_+$ 由 $S_1$ 生成。证明：对任意 $n \ge 0$，$\cO$-模 $\cO(n)$ 由它的整体截面生成，这些截面属于典范映射
\[
S_n \to \Gamma\bigl(\Proj(S), \cO(n)\bigr)
\]
的像。
\end{exercise}

\subsection{应用：闭子概形与紧合性}\label{sec:26.5}

\begin{theorem}\label{thm:26.5.1}
设 $S$ 为 $\Z_{\ge 0}$-分次交换环。
\begin{enumerate}[label=(\arabic*),leftmargin=2.4em,itemsep=0.4ex]
  \item 对任意分次理想 $I \subseteq S$，同态 $\phi \colon S \to S/I$ 诱导闭浸入
  \begin{equation}\label{eq:26.12}
  \Proj(S/I) \to \Proj(S),
  \end{equation}
  其定义理想由 $\widetilde{I}$ 给出。
  \item 设 $S_+$ 由 $S_1$ 中有限多个元素生成。则任意闭浸入 $i \colon X \to \Proj(S)$ 都形如 (26.12)，其中
  \[
  I := \ker\bigl(S \to \Gamma(\cO)_{\mathrm{gr}} \to \Gamma(i_*\cO_X)_{\mathrm{gr}}\bigr).
  \]
\end{enumerate}
\end{theorem}

\begin{proof}
对 (1)，注意 $\phi(S_+)$ 生成 $(S/I)_+$。这说明 $U_+(\phi) = \Proj(S/I)$。因此有定义良好的态射 $\Proj(S/I) \to \Proj(S)$。按定义，在标准开子集 $U_+(f) \subseteq \Proj(S)$ 上，该态射由环同态 $S_{(f)} \to (S/I)_{(f)}$ 的 $\Spec(-)$ 给出，而后者是核为 $I_{(f)}$ 的满射。这给出 (1) 中的断言。

对 (2)，令 $I := \ker(\cO \to i_*\cO_X)$ 为 $i \colon X \to \Proj(S)$ 的定义理想。由定理 \ref{thm:26.4.6}，有
\[
(\Gamma(i_*\cO_X)_{\mathrm{gr}})\widetilde{\ } \xrightarrow{\ \sim\ } i_*\cO_X.
\]
回忆还有 $\widetilde{S} \simeq \cO$。于是
\[
I \simeq \ker\bigl(\widetilde{S} \to (\Gamma(i_*\cO_X)_{\mathrm{gr}})\widetilde{\ }\bigr).
\]
展开定义，右端的态射是对合成
\[
S \to \Gamma(\cO)_{\mathrm{gr}} \to \Gamma(i_*\cO_X)_{\mathrm{gr}}
\]
施行 $(-)\widetilde{\ }$ 得到的。因 $(-)\widetilde{\ }$ 正合，得到 $I \simeq \widetilde{I}$。于是由 (1) 得 $\Proj(S/I) \simeq X$。
\end{proof}

\begin{exercise}\label{exc:26.5.2}
设 $S$ 为 $\Z_{\ge 0}$-分次交换环，且 $S_+$ 由 $S_1$ 生成。设 $I \subseteq S$ 为分次理想，$i \colon \Proj(S/I) \to \Proj(S)$ 为对应的闭浸入。则
\[
i^*\bigl(\cO_{\Proj(S)}(n)\bigr) \simeq \cO_{\Proj(S/I)}(n).
\]
\end{exercise}

\begin{corollary}\label{cor:26.5.3}
设 $S_+$ 由 $S_1$ 中有限多个元素生成。则 $p \colon \Proj(S) \to \Spec(S_0)$ 紧合。
\end{corollary}

\begin{proof}
由引理 \ref{lem:25.4.6} 与引理 \ref{lem:25.4.3}，只需再证 $p$ 泛闭。由引理 \ref{lem:25.3.5}，只需证 $p$ 闭。

设 $X \subseteq \Proj(S)$ 为闭子集，$X_{\mathrm{red}}$ 为对应的既约闭子概形。由定理 \ref{thm:26.5.1}，有 $X \simeq \Proj(S/I)$，其中 $I \subseteq S$ 为某个分次理想。于是可用 $S/I$ 替换 $S$，从而归约为证明 $p(\Proj(S))$ 是 $\Spec(S_0)$ 的闭子集。

按定义，点 $x \in \Spec(S_0)$ 不在 $p$ 的像中，当且仅当 $p^{-1}(x) = \varnothing$。有
\[
\Proj(S) \times_{\Spec(S_0)} \Spec(\kappa_x) \simeq \Proj(S \otimes_{S_0} \kappa_x).
\]
故 $x \notin p(\Proj(S))$ 当且仅当 $\cO_{S \otimes_{S_0} \kappa_x} = 0$。由习题 \ref{exc:26.4.3}，这等价于 $S \otimes_{S_0} \kappa_x$ 最终为零，即当 $n \gg 0$ 时 $S_n \otimes_{S_0} \kappa_x = 0$。注意 $S_n$ 是有限生成 $S_0$-模。故由 Nakayama 引理，上述条件等价于当 $n \gg 0$ 时 $(S_n)_{\mathfrak p} = 0$，其中 $\mathfrak p_x \subseteq S_0$ 是对应于 $x$ 的素理想，即
\[
\Ann_{S_0}(S_n) \not\subseteq \mathfrak p_x, \qquad n \gg 0.
\]
因 $S_+$ 由 $S_1$ 生成，有 $\Ann_{S_0}(S_n) \subseteq \Ann_{S_0}(S_{n+1})$。令 $\mathfrak a$ 为这些零化子的并。于是 $x \notin \im(p)$ 当且仅当 $\mathfrak a \not\subseteq \mathfrak p_x$。换言之，$\im(p) = Z(\mathfrak a)$ 是闭子集，正如所愿。
\end{proof}

\begin{yizhu}
定理 \ref{thm:26.5.1}(2) 之所以取核而不直接取 $I$，是因为「闭子概形 $\to$ 分次理想」并非单射：若 $I$ 与 $I'$ 只在充分高的次数上不同，它们给出同一个闭子概形（比较引理 \ref{lem:26.1.7}）。取核得到的 $I$ 是所有这些分次理想中最大的一个，即把定义理想「饱和化」，对应才成为双射。
[依据：原书 p.131 定理 26.5.1 及其证明；引理 \ref{lem:26.1.7} 说明高次部分的差异在 $\Proj$ 上不可见。]
\end{yizhu}

\begin{yizhu}
紧合性证明的三段式值得记住：有限型（引理 \ref{lem:25.4.6}）、分离（引理 \ref{lem:25.4.3}）已备，故只剩泛闭；而泛闭经引理 \ref{lem:25.3.5}（基变换后仍是 $\Proj$）归约为闭。这说明引理 \ref{lem:25.3.5} 的笛卡儿方块不只是一条形式性质，它把「对任意基变换验证」这一难点整个消掉了。
[依据：原书 p.131 推论 26.5.3 及其证明；引理 \ref{lem:25.3.5} 见原书 p.122。]
\end{yizhu}
